\documentclass[preprint,12pt]{elsarticle}

\usepackage{amsmath,amssymb}
\usepackage{graphicx}
\usepackage{booktabs}
\usepackage{multirow}
\usepackage{array}
\usepackage{adjustbox}
\usepackage{microtype}
\emergencystretch=2em
\usepackage[hidelinks]{hyperref}
\usepackage{xcolor}

\newcommand{\sr}{sim$\to$real}
\newcommand{\Sr}{Sim$\to$real}
\newcommand{\pmsd}[2]{#1\,$\pm$\,#2}
\graphicspath{{figures/}}

\journal{Computers in Biology and Medicine}

\begin{document}

\begin{frontmatter}

\title{Learning panoramic anatomy from CBCT: physics-based simulation and domain randomization for zero- and few-shot segmentation of dental panoramic radiographs}

\author[usd]{Dipesh Tripathi\corref{cor1}}
\ead{dipesh.tripathi@coyotes.usd.edu}
\cortext[cor1]{Corresponding author.}
\affiliation[usd]{organization={University of South Dakota},
            city={Vermillion},
            state={SD},
            postcode={57069},
            country={USA}}

\begin{abstract}
Panoramic radiographs are the most common dental imaging examination, but pixel-level annotation is costly. Because a
panoramic is a projection, a 2D label also cannot show whether two overlapping structures touch in 3D. We investigate how
well segmentation models trained only on panoramics simulated from labelled cone-beam CT (CBCT) transfer to real ones. We built a physics-based simulator of rotational panoramic imaging (moving narrow beam, focal trough,
out-of-trough blur, ghosts, spine shadow), applied it to 532 ToothFairy3 volumes and projected
their labels into overlapping masks for 32 teeth, pulp, the inferior alveolar canal (IAC), incisive canals, maxillary
sinus and restorations, with domain randomization of geometry, metal, contrast, noise and vendor-style processing.
Trained on renders alone, a U-Net segments teeth on 187 real DENTEX panoramics with a Dice of 0.755\,$\pm$\,0.020 and
numbers 90\% of detected teeth correctly; without randomization, transfer fails (Dice 0.28). With four labelled real
images, numbering accuracy reaches 0.933, against 0.454 for a tuned real-only model. Real-only training needs about ten
times more labels to catch up; the advantage disappears at about 100 images. A paired per-factor ablation shows that
image noise and processing govern transfer, while the panoramic geometry factors improve tooth identity.
Physics-based randomization outperforms CycleGAN translation, which altered the anatomy. Finally, a head trained on
third-molar--IAC contact labels computed from CBCT detects contact on real CONTACT crops with an AUC of 0.596, where the
2D overlap is at chance; as pretraining, it helps up to about 40 labelled real crops.
\end{abstract}

\begin{highlights}
\item Physics-based panoramic simulator projects overlapping labels from labelled CBCT
\item Models trained only on simulated panoramics segment and number real teeth
\item Four real labels plus simulation match real-only training on about 40 images
\item Paired ablation: image noise drives transfer; geometry helps tooth identity
\item CBCT-derived labels detect third-molar--canal contact where 2D overlap fails
\end{highlights}

\begin{keyword}
Panoramic radiography \sep Cone-beam computed tomography \sep Synthetic data \sep Domain randomization \sep
Sim-to-real transfer \sep Tooth segmentation \sep Mandibular canal
\end{keyword}

\end{frontmatter}

% =====================================================================================================================
\section{Introduction}

The panoramic radiograph, or orthopantomogram, is the standard first-line image in dentistry. A single low-dose exposure
shows both dental arches, the mandible, the inferior alveolar canals (IAC) and the lower maxillary sinuses. It is used to
plan extractions and implants, find caries and periapical disease, and assess the relation between an impacted third
molar and the IAC before surgery~\cite{mallya2018}. Deep learning on panoramics has advanced quickly for tooth detection
and numbering~\cite{tuzoff2019,jader2018}, diagnosis~\cite{hamamci2023}, and segmentation of the canal and other
structures~\cite{vinayahalingam2019,cha2021,aung2024}. But these methods depend on pixel-level labels drawn by dentists
on 2D images. Such labels are expensive and inconsistent between datasets. For thin or low-contrast structures such as
the canal, pulp and sinus floor, they are also hard to draw at all. Public panoramic datasets therefore mostly label
teeth. We found no public expert segmentations of the sinus, pulp or incisive canal on panoramics.

A panoramic image also has a problem that annotation cannot fix. It is a curved, blurred projection, so structures at
different depths overlap. Whether a third-molar root actually touches the IAC is the question that decides the risk of
nerve injury during extraction, and it cannot be read reliably from the panoramic itself. The classical radiographic
signs of Rood and Shehab~\cite{rood1990} are only moderately predictive. Studies that use CBCT as the reference standard
report that both clinicians and convolutional networks find the task hard on
panoramics~\cite{fukuda2020,jeon2023,ulusoy2025}. A 2D label drawn on the panoramic cannot carry the 3D truth.

Cone-beam CT offers both things the panoramic lacks. Public CBCT datasets now come with dense voxel labels: ToothFairy3
provides 532 volumes with 77 classes, including every tooth, its pulp, both IACs and the incisive
canals~\cite{toothfairy3,cipriano2022}. If a realistic panoramic can be simulated from such a volume, every label can be
projected exactly, overlaps included, and 3D relations such as root--canal contact can be computed directly. This is the
logic of in-silico training. SyntheX showed that realistic simulation from CT can replace real training data for X-ray
analysis~\cite{gao2023}. Domain randomization makes models robust to the remaining gap by varying the simulator
widely~\cite{tobin2017,billot2023}. In dentistry, simulated panoramics from CBCT have so far been used mainly as an
intermediate for 3D reconstruction~\cite{song2021,park2024,szczepanski2026,yue2026}, or as pretraining data for tooth
segmentation~\cite{wang2025}. How far a panoramic simulator alone can carry multi-structure segmentation and 3D-relation
prediction on real images, and which parts of the simulation matter, has not been tested.

This paper addresses that question. Our contributions are:
\begin{enumerate}
\item \textbf{A physics-based panoramic simulator over ToothFairy3.} It models the rotational narrow-beam acquisition,
including the focal trough, blur outside it, ghosts and the spine shadow. It projects overlapping multi-label masks for
32 teeth by FDI number, pulp, both IACs, incisive canals, the maxillary sinus, crowns, bridges and implants. Public real
labels exist only for the teeth and the canal, so these two are validated on real images; the other structures are
evaluated on held-out synthetic data.
\item \textbf{Panoramic-specific domain randomization with a paired per-factor ablation.} Nine factor groups are
switched off one at a time with all other random draws held fixed, over three seeds. Follow-up experiments retune the
most influential ranges.
\item \textbf{A sim-to-real evaluation on real panoramics.} It covers zero-shot transfer, few-shot fine-tuning against
real-only training at 1--100\% of the labels, a no-randomization baseline, and a CycleGAN image-translation baseline in
the style of NeBLa~\cite{park2024}.
\item \textbf{A 3D-relation head.} It predicts true third-molar--IAC contact from the 2D image alone, trained with
contact labels computed automatically from each CBCT and tested on real CONTACT crops with CBCT-confirmed labels.
\end{enumerate}

The main findings are as follows. First, a model trained on simulation alone segments and numbers teeth on real
panoramics. With a handful of real labels it matches real-only training that uses about ten times more, although the
advantage disappears once about 100 labelled images are available. Randomization is what makes this work: unrandomized
simulation does not transfer. Second, image-level realism (noise and vendor processing) dominates transfer, while
geometry helps tooth identity. Third, physics-based randomization preserves label correctness, whereas unpaired
translation does not. Fourth, a contact signal learned from CBCT transfers weakly to real images, while 2D overlap
carries none; it is a useful starting point only when few real contact labels exist.

% =====================================================================================================================
\section{Related work}

\subsection{Deep learning on panoramic radiographs}
Early work detected and numbered teeth with two-stage detectors~\cite{tuzoff2019} and segmented them with Mask
R-CNN~\cite{jader2018}, often on private data. Public benchmarks followed. The Tufts Dental Database provides 1,000
panoramics with expert tooth and abnormality labels~\cite{panetta2022}. DENTEX, a MICCAI 2023 challenge, provides
hierarchically labelled quadrants, FDI enumeration and diagnoses, plus a large unlabelled pool~\cite{hamamci2023}.
Non-tooth anatomy has received less attention and smaller datasets. Vinayahalingam et al.~\cite{vinayahalingam2019}
segmented third molars and the IAC on 81 panoramics. Cha et al.~\cite{cha2021} performed panoptic segmentation of the
sinus, jaws, canal, teeth and implants on 51 images. Aung et al.~\cite{aung2024} studied canal segmentation across
devices. All of these rely on 2D expert labels drawn on the radiograph. Our work instead derives labels from CBCT and
needs no 2D annotation for training.

\subsection{Third-molar--canal relation}
Rood and Shehab~\cite{rood1990} described radiographic signs on panoramics, such as darkening of the root, interruption
of the canal's white line and diversion of the canal, that are associated with nerve injury. CBCT is the reference
standard for actual contact. Several deep-learning studies classify the relation directly from panoramic crops using
CBCT-derived labels~\cite{fukuda2020,jeon2023}, with accuracies around 0.7--0.8 on held-out data. The CONTACT
dataset~\cite{ulusoy2025} is the first public panoramic dataset of this kind. It has 1,478 third-molar--canal pairs from
1,010 patients, semantic masks, expert-scored radiographic features and CBCT-confirmed contact labels. Its authors report
that learned models and clinicians alike find the ``overlap without contact'' case hardest. All prior approaches learn
contact from real labelled panoramics. We learn it from CBCT geometry alone and use CONTACT only for testing and few-shot
experiments.

\subsection{Simulated X-rays and sim-to-real transfer}
Digitally reconstructed radiographs (DRRs) from CT have long been used in registration. DeepDRR added realistic material
decomposition, scatter and noise, and showed that models trained on such DRRs transfer to real fluoroscopy, while models
trained on naive DRRs do not~\cite{unberath2018,unberath2019}. DiffDRR made DRR rendering
differentiable~\cite{gopalakrishnan2022}. SyntheX scaled this idea and found that realistic simulation combined with
strong domain randomization can match or exceed real-data training for hip imaging, surgical tools and COVID-19
lesions~\cite{gao2023}. AnyCXR applies multi-stage randomization to chest X-ray anatomy at arbitrary projection
angles~\cite{dong2025}. Domain randomization itself originates in robotics~\cite{tobin2017}. In medical imaging, SynthSeg
showed that fully randomizing contrast and resolution yields brain-MRI segmentation that works on any
contrast~\cite{billot2023}. Panoramic imaging differs from these settings because the image is not a single projection.
A narrow beam rotates around the patient while the detector moves, so a curved layer is sharp and everything else is
smeared. Most of what makes a panoramic look like a panoramic comes from this geometry. To our knowledge, the effect of
randomizing it has not been measured.

\subsection{Simulated panoramics from CBCT}
Simulated panoramics have mostly served 3D reconstruction from a single panoramic. Oral-3D~\cite{song2021} and
X2Teeth~\cite{liang2021} reconstruct bone or teeth from CBCT-derived panoramics. PX2Tooth~\cite{ma2024} reconstructs
tooth point clouds and builds paired data by resampling CBCT along the dental arch. NeBLa~\cite{park2024} simulates
panoramics with Beer--Lambert attenuation along rays following panoramic rotation (SimPX). At test time it translates
real panoramics into the simulated style with a learned translation module. X-Splat~\cite{szczepanski2026} and
HyDAR-Pano3D~\cite{yue2026} train on simulated panoramic--CBCT pairs; X-Splat uses SimPX-style rendering on ToothFairy
data. For 2D analysis, PXseg~\cite{wang2025} pretrains tooth segmentation and numbering on CBCT-derived panoramics and
then trains on real labelled images, reporting a 4.6\% gain in tooth Dice over no pretraining. These works either target
teeth only, depend on real labels or a translation step at test time, or do not model the panoramic blur and ghosting.
Our no-randomization baseline (Section~\ref{sec:setup}) approximates PXseg-style pretraining within our pipeline. None
studies which simulation factors matter. Our simulator models the rotational time-delay integration geometry explicitly.
We evaluate zero-shot transfer for many structures at once and ablate each factor.

\subsection{Unpaired image translation}
An alternative to making the simulator realistic is to make its images look real afterwards, with unpaired translation
such as CycleGAN~\cite{zhu2017}. Translation can close the appearance gap, but distribution-matching losses are known to
add or remove image content when the content statistics of the two domains differ~\cite{cohen2018}. Labels attached to
the source image then no longer match the translated one. We compare against this approach directly.

% =====================================================================================================================
\section{Methods}

\subsection{Data}
\textbf{ToothFairy3}~\cite{toothfairy3} contains 532 CBCT volumes at 0.3\,mm isotropic resolution in Hounsfield units
(HU), with voxel labels for 77 classes. These include the jaws, every tooth by FDI number, each tooth's pulp, both IACs,
the incisive and lingual canals, the maxillary sinuses, and crowns, bridges and implants. The volumes come in three
sets: A (417 volumes, mostly mandible-centred field of view), B (63, full face) and C (52, a different scanner). The
voxel data are stored in RPI order although the header states LPS. We verified this with anatomical tests on all 532
label volumes and correct it in memory. The officially distributed reorientation script yields an upside-down,
front-to-back mirrored volume, so we did not use it. Upper-jaw structures are present mainly in Set B: the maxillary
sinus appears in about 50 volumes and the upper jawbone in 43\% of them. Each volume carries per-channel ``annotated''
flags and a field-of-view mask, so the loss ignores structures that a volume does not contain. We split the volumes once
into 374 training, 53 validation and 105 test volumes, stratified by set, metal and third-molar presence.

\textbf{DENTEX}~\cite{hamamci2023} supplies real panoramics with FDI-enumerated tooth polygons, plus 2,098 unlabelled
images. The same radiograph appears under different file names across and within its subsets. We grouped exact
duplicates with a 2048-bit difference hash and made all splits over duplicate groups. This gives a few-shot training pool
of 384 labelled images, 62 validation images and a test set. Evaluation uses one image per duplicate group among the test
images that have tooth polygons: 187 images.

\textbf{CONTACT}~\cite{ulusoy2025} provides 1,478 third-molar region crops from 1,010 patients, at $227\times227$
pixels, with molar and canal masks, six expert-scored radiographic signs and a CBCT-confirmed contact label. Each crop
belongs to one of three groups: (1) overlap with contact, (2) no overlap and no contact, (3) overlap without contact. We
keep the official test set of 444 crops (330 / 34 / 80 by group) and hold out 15\% of the official training set as
validation (878 training, 156 validation crops). The release has no patient identifiers, so the training and validation
crops may share patients. The test set is unaffected.

All datasets are used under their non-commercial research licences.

\subsection{Rotational panoramic simulator}
A panoramic unit rotates a narrow vertical slit beam and a detector around the patient. The detector moves past the slit
at a speed matched to the beam's sweep along the arch. Points on a curved focal layer, the focal trough, project to the
same detector location throughout the exposure and are sharp. Points in front of or behind it drift across the detector
and are blurred horizontally. Structures on the far side of the head, such as the opposite ramus and the spine, are
smeared into ghost images. We simulate this time-delay integration (TDI) process directly (Fig.~\ref{fig:sim}).

\textbf{Arch curve.} For each volume we fit the focal-trough centre curve in polar form $\rho(\theta)$ around a centre
inside the mandibular U. At every $2^\circ$ direction, the radius is the middle of the combined 5th--95th-percentile
radial range of the teeth and canal voxels; where neither exists, we use the mandibular body. The curve is smoothed over
$\theta$ (Gaussian, $\sigma = 6^\circ$) and extended 10\,mm at each end. An earlier fit to the teeth alone ended too
early in partially edentulous jaws and cut off the posterior canal in 125 of 532 volumes. The combined fit keeps 100\% of
the canal within $\pm12$\,mm of the curve in 531 volumes, and at least 99.3\% of the tooth voxels in every volume. For
edentulous volumes, the arch is fitted to the alveolar ridge.

\textbf{Beam geometry.} Time is parametrised by arc length $s$ along the arch, in steps of 0.3\,mm. At each step the
central ray passes through the arch point $A(s)$, tilted upward by $5^\circ$. The source lies 400\,mm behind the trough
and the detector 100\,mm in front, for a magnification of $M = 1.25$. A 6\,mm slit is modelled by casting one ray per
slit pixel. The rotation centre for each step is set from the local radius of curvature of the arch, clipped to
[15, 80]\,mm and pulled toward the arch normal. This keeps the beams nearly perpendicular to the arch at the incisors and
premolars, and oblique through the molar region as in real units. The detector advances by $M\cos\delta$ per step, where
$\delta$ is the beam obliquity, so that points on the trough stay fixed on the detector. The result is a narrow trough
anteriorly and a wide one posteriorly. We checked this with a bead phantom: beads inside the trough render 0.9--1.2\,mm
wide at all heights, off-trough beads blur and drift as expected, and the far side produces ghosts.

\textbf{Attenuation and image formation.} Volumes are converted to linear attenuation
$\mu = \mu_\mathrm{water}(1 + \mathrm{HU}/1000)$, clipped at zero, with $\mu_\mathrm{water} = 0.02$\,mm$^{-1}$
nominally. Rays are sampled through the volume with trilinear interpolation on the GPU. The detector records the
transmission $T = \exp(-\int \mu\,\mathrm{d}s)$, averaged over the exposure. The image is $-\log T$, and the raw mean
transmission is stored so that noise and processing can be applied later. Rendering takes about 5\,s per image on a
V100 GPU.

\textbf{Label projection.} For each detector pixel and channel, we record whether the ray hits the structure within a
depth window around the focal point, for at least 50\% of the exposure. Using the depth window rather than the whole ray
keeps far-side structures such as the opposite canal out of the masks. Channels are projected independently, so overlaps
are preserved: a tooth and the canal it overlaps are both labelled at the same pixel. We project 43 channels: 11
structure channels (teeth union, pulp, left and right IAC, incisive canals, sinus, crown, bridge, implant, mandible,
maxilla) and 32 FDI tooth channels. Panoramics are cropped vertically from the top of the field of view to the lower
border of the mandible plus 8\,mm, which removes most of the CBCT chin rest. The images are resampled to a
$256\times640$ canvas at 0.3\,mm per pixel.

\begin{figure}[t]
\centering
\includegraphics[width=\textwidth]{fig_simulator_DATA.png}
\caption{The simulator on one ToothFairy3 volume (Set B). (a) Nominal render. (b) The same render with its projected
labels; channels are projected independently, so teeth, pulp, canals and sinus overlap where they do in projection.
(c) Randomized training renders of the same volume: geometry and volume factors are fixed when a render is made, and the
image-level factors are drawn as in training.}
\label{fig:sim}
\end{figure}

\subsection{Domain randomization}
Every randomization factor has a nominal value, used for the validation and test renders, and a sampling range for
training. The factors act at three levels (Table~\ref{tab:dr}). \emph{Volume factors} change the attenuation before
projection: tissue contrast, spine density, and metal density (restoration labels scaled by a factor). \emph{Geometry
factors} change the acquisition: trough width (a scale on the rotation radius), slit width (blur outside the trough),
ghost weight (the far-side contribution), and arch-curve jitter (shift, rotation, scale and head tilt). \emph{Image
factors} are applied on the fly during training: Poisson photon noise, Gaussian read noise and detector blur, and
vendor-style processing (large-kernel background subtraction, unsharp masking, gamma).

\begin{table}[t]
\centering\small
\caption{Randomization factors. Geometry and volume factors are fixed when a render is made; image factors are drawn
again for every training sample. U = uniform, log-U = log-uniform.}
\label{tab:dr}
\begin{adjustbox}{max width=\linewidth}
\begin{tabular}{@{}llp{9.2cm}@{}}
\toprule
Factor & Level & Parameters (nominal $\to$ training range) \\
\midrule
Trough width & geometry & rotation-radius scale 1 $\to$ log-U(0.7, 1.6) \\
Blur outside trough & geometry & slit width 6 $\to$ U(3, 10)\,mm \\
Ghosts & geometry & far-side weight 1 $\to$ U(0.4, 1.3) \\
Arch jitter & geometry & AP shift $\pm$4\,mm, lateral $\pm$2\,mm, rotation $\pm4^\circ$, scale $\pm$5\%, tilt
$5^\circ \to$ U($0^\circ$, $10^\circ$) \\
Spine & volume & spine attenuation scale 1 $\to$ U(0.3, 1.6) \\
Metal & volume & restoration attenuation scale 1 $\to$ log-U(1, 3) \\
Contrast & volume & $\mu_\mathrm{water}$ 0.02 $\to$ U(0.016, 0.026)\,mm$^{-1}$; bone contrast 1 $\to$ U(0.75, 1.3) \\
Noise & image & photons $\infty \to$ log-U($2\cdot10^3$, $5\cdot10^4$); read noise $\leq$\,1\%; PSF
$\sigma \leq 0.4$\,mm \\
Processing & image & background $\sigma$ 12 $\to$ U(6, 25)\,mm, amount 0.85 $\to$ U(0.6, 0.95); unsharp $\leq$\,1;
$\gamma \to$ log-U(0.7, 1.4) \\
\bottomrule
\end{tabular}
\end{adjustbox}
\end{table}

For each training volume we render a nominal image (v0) and eight randomized variants (v1--v8), giving 3,843 training and
validation renders. Random draws come from a generator seeded per volume and variant, and every sampler is always drawn,
even when its factor is switched off. Switching a factor off therefore changes only that factor; all other parameters of
the render stay identical. This allows a \emph{paired} ablation: each ablated render set differs from the reference only
in the removed factor.

\subsection{Segmentation model and training}
We use a U-Net~\cite{ronneberger2015} with a ResNet-34 encoder~\cite{he2016} initialised from ImageNet, one input channel
and 43 sigmoid outputs. Channels overlap, so each is a separate binary problem. For the same reason we did not use
nnU-Net~\cite{isensee2021}, which assumes one label per pixel. The loss is per-pixel binary cross-entropy plus a
per-image, per-channel soft Dice term. Both are restricted to pixels where the channel's label is valid, according to the
field of view and the annotated flags, and Dice applies only where the structure is present. On real DENTEX images, only
the teeth union and the 32 FDI channels are supervised.

Augmentation consists of a random crop of 80--100\% of the image and a horizontal flip. The flip swaps left and right
channels: the canals, and FDI quadrants 1$\leftrightarrow$2 and 3$\leftrightarrow$4. Real images are framed with one
fixed relative crop, derived from tooth extents in the few-shot training pool, so that they match the synthetic canvas.
We train with AdamW (learning rate $10^{-3}$, weight decay $10^{-4}$), a one-cycle schedule, batch size 16 and mixed
precision, for 60 epochs on one V100 GPU (about 2.3 hours). For the zero-shot model, the checkpoint is selected on the
mean Dice over synthetic validation renders only. Real images are never used for training or selection.

\subsection{Third-molar crops and the contact head}
CONTACT provides crops, not whole panoramics, so we train a crop model on matching synthetic crops. From the projected
masks we locate each lower third molar (FDI 38 or 48) and the canal on the same side, and cut a square crop of 75 pixels
(about 22.5\,mm) around the molar--canal junction. This size matches the third-molar extent in CONTACT, about
$172\times140$ pixels of 227. During training the crop scale is jittered by $\times$0.8--1.25 and the position by
$\pm$12\%, with random flips. The model has two channels, molar and canal, and is initialised from the panoramic model:
its molar head is the mean of the tooth-38 and tooth-48 weights, and its canal head the mean of the two canal heads.
Synthetic crops come from 0.3\,mm renders and are softer than CONTACT's crops. For evaluation we therefore downsample each
real crop to the synthetic resolution and resize it back (matched resolution).

\textbf{3D contact labels.} For every lower third molar in ToothFairy3 we compute the minimum Euclidean distance between
the tooth and the canal label on the same side, using a Euclidean distance transform with the true voxel spacing within
a 15\,mm margin. We define contact a priori as the labels touching, including diagonally: a distance of at most
0.52\,mm, or $\sqrt{3}$ times the 0.3\,mm voxel size. Of 583 lower third molars, 195 (33\%) are in contact: 411 / 58 /
114 molars in the training, validation and test volumes. No manual labelling is involved.

\textbf{Contact head.} A classification head on the encoder bottleneck (global pooling, dropout 0.2, one logit) is
trained with binary cross-entropy on these labels, alongside the segmentation loss (weight 0.5). The model is initialised
from the crop model and selected on synthetic validation contact AUC. As a no-learning baseline, we also read contact
from the predicted 2D molar--canal overlap fraction, which is the information a panoramic shows directly.

% =====================================================================================================================
\section{Experimental setup}
\label{sec:setup}

\textbf{Zero-shot.} The simulation-only model is evaluated on the 187 DENTEX test images and on nominal renders of the
105 held-out ToothFairy3 test volumes.

\textbf{No-randomization baseline.} To isolate the effect of randomization, we train the same model on the nominal
renders only (v0 of each training volume, no image-level randomization), for the same number of optimisation steps as
the main model and with the same selection rule, using three seeds. This is the closest analogue within our pipeline to
pretraining on unrandomized CBCT-derived panoramics, as in PXseg~\cite{wang2025}.

\textbf{Few-shot.} Each arm uses 1, 5, 10, 25, 50 or 100\% of the 384-image pool (4 to 384 images), with three seeds that
change the image subset and the training run, and selects checkpoints on the 62 real validation images. \emph{\Sr}
fine-tunes the zero-shot model (seed 0) with learning rate $3\cdot10^{-4}$ for 40 epochs. \emph{No-randomization
\sr} fine-tunes the no-randomization model (seed 0) with the same schedule. \emph{Real-only} trains the same network from
ImageNet initialisation. A model trained from ImageNet needs a longer schedule than one being fine-tuned: with the
fine-tuning schedule, real-only per-tooth Dice was still rising at the last epoch. We therefore tuned the real-only arm
separately (learning rate $10^{-3}$, 120 epochs, validation every third epoch), and report only this tuned baseline. To
check that the comparison does not favour real-only by giving it three times the training, we also trained the \sr\ arm
for 120 epochs. FDI Dice changed by at most 0.006 and numbering accuracy by at most 0.002 at any budget, and the selected
checkpoints came early in both schedules (Table~\ref{tab:sched}). We therefore report the 40-epoch schedule.

\textbf{Paired ablation.} The reference uses variants v1--v4 of every training volume with all factors on, trained for
40 epochs. Each ablation run replaces these with the paired render set in which one factor is at its nominal value; for
the image-level factors, the factor is simply not applied. Each configuration is trained with three seeds and compared
with the reference trained with the same seed.

\textbf{CycleGAN baseline.} Following the NeBLa strategy of bridging appearance by learned translation, we train a
standard CycleGAN~\cite{zhu2017} (ResNet-9 generator, $70\times70$ PatchGAN, least-squares loss, cycle weight 10,
identity weight 0.5) for 40,000 iterations between nominal synthetic renders and the 2,098 unlabelled DENTEX images, none
of which shares a duplicate group with the test set. We then train the segmentation model, with the zero-shot schedule,
on synthetic renders passed through the frozen generator together with their original masks. We do this both without
and with our image-level randomization.

\textbf{CONTACT.} Zero-shot crop and contact models are tested on the 444 official test crops. For few-shot contact, we
compare fine-tuning the contact model with training from ImageNet on 1--100\% of the 878 training crops (masks and
contact labels), with three seeds and selection on validation contact AUC. Both arms use the same schedule (learning rate
$10^{-3}$, 90 epochs, validation every third epoch), so they differ only in initialisation. We first ran both with the
shorter fine-tuning schedule (learning rate $3\cdot10^{-4}$, 30 epochs). That undertrained the real-only arm (canal Dice
0.03 to 0.14 lower, contact AUC 0.05 lower with all crops). For \sr, the longer schedule mattered only together with the
higher learning rate (Table~\ref{tab:sched}). On validation, the three \sr\ schedules were indistinguishable (contact AUC
within 0.013 at every budget), so we fixed the schedule by matching the real-only arm rather than by test performance.

\textbf{Metrics.} On DENTEX we report Dice of the teeth union and mean Dice over the FDI channels present in the ground
truth. DENTEX polygons are coarse: straight-edged, and including the gaps between teeth. We therefore also report three
boundary-insensitive tooth-level measures. A ground-truth tooth counts as \emph{detected} if at least 50\% of its polygon
is covered by predicted tooth pixels. \emph{Numbering accuracy} is the fraction of detected teeth whose largest-overlap
predicted FDI channel is correct. \emph{Presence F1} measures, per FDI number, whether a tooth is predicted (area
$\geq$\,300 pixels) when present. On CONTACT we report molar and canal Dice, and the AUC of contact (group 1 vs.\ groups
2 + 3), overall and for the hard comparison of group 1 vs.\ 3. Confidence intervals are 95\% percentile bootstraps over
test images (2,000 resamples).

% =====================================================================================================================
\section{Results}

\subsection{Zero-shot transfer}
Trained only on simulated panoramics, the model segments the teeth on the 187 real DENTEX test images with a Dice of
\pmsd{0.755}{0.020} over three training seeds (Table~\ref{tab:zs}). It also numbers them: mean FDI Dice is
\pmsd{0.582}{0.019}, 87.1\% of annotated teeth are detected, and \pmsd{90.0}{0.1}\% of detected teeth receive the
correct FDI number. Numbering accuracy is almost identical across seeds, while Dice and detection vary more. The model
never saw a real image or a real label.

Visual inspection explains much of the remaining Dice gap (Fig.~\ref{fig:dentex}). The predictions follow tooth contours
and leave interproximal gaps empty, whereas the DENTEX polygons are straight-edged and include those gaps. Dice therefore
partly measures agreement with an annotation style, which is why we also report the tooth-level measures. On the 105
held-out synthetic test volumes the model reaches a teeth Dice of 0.882, IAC 0.745 / 0.763 (left / right), sinus 0.709,
pulp 0.668 and incisive canal 0.460. On real images the predicted canals are anatomically plausible: they run below the
molar roots and rise into the ramus. Real canal accuracy is measured on CONTACT in Section~\ref{sec:contact}.

All downstream experiments (few-shot, crops, contact) start from training seed 0, which happens to be the best of the
three seeds on DENTEX. We report its numbers with bootstrap confidence intervals alongside the mean, so that the starting
point of those experiments is clear.

\begin{table}[t]
\centering\small
\caption{Zero-shot results of the simulation-only model on the DENTEX test set (187 images). Seed 0 initialises every
downstream experiment.}
\label{tab:zs}
\begin{adjustbox}{max width=\linewidth}
\begin{tabular}{@{}lcc@{}}
\toprule
Measure & Mean $\pm$ sd, 3 seeds & Seed 0 [95\% CI] \\
\midrule
Teeth Dice & \pmsd{0.755}{0.020} & 0.774 [0.765, 0.783] \\
FDI Dice (mean over teeth present) & \pmsd{0.582}{0.019} & 0.604 [0.582, 0.626] \\
Tooth detection rate & \pmsd{0.871}{0.038} & 0.912 [0.895, 0.928] \\
Numbering accuracy & \pmsd{0.900}{0.001} & 0.899 [0.880, 0.918] \\
Presence F1 & \pmsd{0.929}{0.013} & 0.943 \\
\bottomrule
\end{tabular}
\end{adjustbox}
\end{table}

\begin{figure}[t]
\centering
\includegraphics[width=\textwidth]{fig_dentex_overlays_DATA_paper.png}
\caption{Zero-shot predictions of the simulation-only model on real DENTEX test panoramics~\cite{hamamci2023}. The model
has never seen a real image or label. Predicted tooth outlines follow the tooth contours, while the ground-truth polygons
are straight-edged and include interproximal gaps. Canals and sinus are predicted but have no ground truth in DENTEX.}
\label{fig:dentex}
\end{figure}

\subsection{Few-shot: simulation replaces most labels when labels are scarce}
Simulation pretraining helps most when real labels are scarce (Table~\ref{tab:fewdentex}, Fig.~\ref{fig:fewdentex}).
With four labelled images, \sr\ reaches an FDI Dice of \pmsd{0.702}{0.006} and a numbering accuracy of
\pmsd{0.933}{0.006}. The tuned real-only model reaches \pmsd{0.169}{0.042} and \pmsd{0.454}{0.073} on the same images.
Real-only training needs about ten times more labels to catch up: with 39 images it reaches an FDI Dice of 0.696 and a
numbering accuracy of 0.923, close to what \sr\ achieves with four. With 20 images \sr\ is still clearly ahead (FDI 0.720
vs.\ 0.620, numbering 0.935 vs.\ 0.869). From about 100 labelled images the two are equivalent; with all 384, real-only
is marginally ahead on FDI Dice (0.779 vs.\ 0.774) and level on numbering (0.952 vs.\ 0.950). Simulation therefore
replaces most real labels in the low-data regime but adds little once a moderately sized labelled set exists.
Teeth-union Dice behaves similarly but with smaller gaps: the outline of the dentition in DENTEX's annotation style is
learned from few images, while the identity of each tooth is not.

Randomization matters here too. Fine-tuning from the no-randomization model also helps at small budgets (FDI 0.619 with
four images), so even unrandomized simulation carries useful anatomical information. But it is below the randomized model
at every budget, and from 39 images it falls behind real-only training (FDI 0.751 vs.\ 0.779 with all 384 images).
Pretraining on unrandomized renders seems to leave the network committed to the simulator's appearance, which then has
to be unlearned.

\begin{figure}[t]
\centering
\includegraphics[width=\textwidth]{fig_fewshot_dentex.pdf}
\caption{Few-shot results on the DENTEX test set as a function of the number of labelled real images (mean over three
seeds; bands show $\pm$ one sd). The dashed line is the zero-shot model.}
\label{fig:fewdentex}
\end{figure}

\begin{table}[t]
\centering\small
\caption{Few-shot results on the DENTEX test set, mean $\pm$ sd over three seeds. RO = real-only (ImageNet
initialisation, tuned schedule); NoDR = fine-tuned from the no-randomization model; DR = fine-tuned from the randomized
zero-shot model. Bold marks a clear best (difference larger than the spread across seeds).}
\label{tab:fewdentex}
\begin{adjustbox}{max width=\linewidth}
\begin{tabular}{@{}rlccc@{}}
\toprule
Labelled images & Arm & Teeth Dice & FDI Dice & Numbering \\
\midrule
\multirow{3}{*}{4 (1\%)} & RO & \pmsd{0.793}{0.009} & \pmsd{0.169}{0.042} & \pmsd{0.454}{0.073} \\
 & NoDR & \pmsd{0.839}{0.006} & \pmsd{0.619}{0.011} & \pmsd{0.886}{0.011} \\
 & DR & \textbf{\pmsd{0.860}{0.013}} & \textbf{\pmsd{0.702}{0.006}} & \textbf{\pmsd{0.933}{0.006}} \\
\midrule
\multirow{3}{*}{20 (5\%)} & RO & \pmsd{0.865}{0.006} & \pmsd{0.620}{0.020} & \pmsd{0.869}{0.016} \\
 & NoDR & \pmsd{0.862}{0.007} & \pmsd{0.662}{0.006} & \pmsd{0.901}{0.003} \\
 & DR & \textbf{\pmsd{0.875}{0.003}} & \textbf{\pmsd{0.720}{0.007}} & \textbf{\pmsd{0.935}{0.002}} \\
\midrule
\multirow{3}{*}{39 (10\%)} & RO & \pmsd{0.881}{0.000} & \pmsd{0.696}{0.006} & \pmsd{0.923}{0.005} \\
 & NoDR & \pmsd{0.872}{0.002} & \pmsd{0.685}{0.007} & \pmsd{0.911}{0.003} \\
 & DR & \pmsd{0.881}{0.004} & \textbf{\pmsd{0.736}{0.005}} & \textbf{\pmsd{0.940}{0.001}} \\
\midrule
\multirow{3}{*}{96 (25\%)} & RO & \pmsd{0.888}{0.001} & \pmsd{0.738}{0.007} & \pmsd{0.941}{0.007} \\
 & NoDR & \pmsd{0.880}{0.001} & \pmsd{0.708}{0.002} & \pmsd{0.920}{0.003} \\
 & DR & \pmsd{0.886}{0.002} & \textbf{\pmsd{0.749}{0.002}} & \pmsd{0.940}{0.002} \\
\midrule
\multirow{3}{*}{192 (50\%)} & RO & \pmsd{0.890}{0.001} & \pmsd{0.756}{0.005} & \pmsd{0.946}{0.003} \\
 & NoDR & \pmsd{0.886}{0.001} & \pmsd{0.734}{0.004} & \pmsd{0.933}{0.002} \\
 & DR & \pmsd{0.890}{0.001} & \pmsd{0.762}{0.002} & \pmsd{0.949}{0.001} \\
\midrule
\multirow{3}{*}{384 (100\%)} & RO & \pmsd{0.894}{0.002} & \pmsd{0.779}{0.002} & \pmsd{0.952}{0.002} \\
 & NoDR & \pmsd{0.890}{0.001} & \pmsd{0.751}{0.004} & \pmsd{0.941}{0.002} \\
 & DR & \pmsd{0.894}{0.000} & \pmsd{0.774}{0.002} & \pmsd{0.950}{0.002} \\
\bottomrule
\end{tabular}
\end{adjustbox}
\end{table}

\subsection{Which randomization factors matter}
The paired ablation (Table~\ref{tab:abl}, Fig.~\ref{fig:abl}) separates the factors into three groups. The reference
reaches a teeth Dice of \pmsd{0.734}{0.004} and a numbering accuracy of \pmsd{0.911}{0.004} with the shorter ablation
schedule.

\emph{Image-level noise is essential.} Without it, teeth Dice drops by 0.436 [$-$0.547, $-$0.344], while accuracy on
the synthetic test set stays the same or improves: without noise the model fits the simulator but does not transfer.
\emph{Vendor-style processing} comes second ($-$0.066 [$-$0.095, $-$0.023]). Both effects appear in every seed.

\emph{The panoramic geometry factors} (trough width, ghosts, blur outside the trough, arch jitter, spine) change teeth
Dice by at most $\pm$0.03. Removing any of them lowers numbering accuracy, by 0.008 to 0.022, and the drop appears in all
three seeds for each of them. These factors help the model tell neighbouring teeth apart rather than outline them.

\emph{Metal randomization, as configured, hurts.} Removing it raises teeth Dice by 0.033 [0.022, 0.047] in every seed,
at a small cost in numbering ($-$0.015).

\begin{table}[t]
\centering\small
\caption{Paired per-factor ablation on the DENTEX test set. Each $\Delta$ is relative to the reference trained with the
same seed; brackets give the minimum and maximum over three seeds.}
\label{tab:abl}
\begin{adjustbox}{max width=\linewidth}
\begin{tabular}{@{}lccc@{}}
\toprule
Factor removed & $\Delta$ teeth Dice & $\Delta$ numbering & Same sign in 3/3 seeds \\
\midrule
Noise & $-$0.436 [$-$0.547, $-$0.344] & $-$0.062 [$-$0.181, +0.026] & teeth \\
Processing & $-$0.066 [$-$0.095, $-$0.023] & $-$0.021 [$-$0.038, +0.006] & teeth \\
Spine & $-$0.011 [$-$0.038, +0.009] & $-$0.008 [$-$0.015, $-$0.001] & numbering \\
Arch jitter & $-$0.009 [$-$0.043, +0.016] & $-$0.018 [$-$0.039, $-$0.000] & numbering \\
Contrast & +0.001 [$-$0.015, +0.021] & $-$0.008 [$-$0.019, +0.005] & -- \\
Trough width & +0.006 [$-$0.004, +0.015] & $-$0.022 [$-$0.028, $-$0.012] & numbering \\
Ghosts & +0.016 [$-$0.010, +0.033] & $-$0.020 [$-$0.031, $-$0.011] & numbering \\
Blur outside trough & +0.026 [+0.013, +0.047] & $-$0.018 [$-$0.026, $-$0.007] & teeth, numbering \\
Metal & +0.033 [+0.022, +0.047] & $-$0.015 [$-$0.027, +0.005] & teeth \\
\bottomrule
\end{tabular}
\end{adjustbox}
\end{table}

\begin{figure}[t]
\centering
\includegraphics[width=\textwidth]{fig_ablation.pdf}
\caption{Paired ablation: change in DENTEX teeth Dice (left) and numbering accuracy (right) when one factor is switched
off. Bars show the mean, dots the three seeds; the noise bar extends beyond the axis.}
\label{fig:abl}
\end{figure}

\textbf{Retuning the ranges.} The ablation suggested two changes, which we tested under the same protocol
(Table~\ref{tab:follow}). Narrowing the metal range from $\times$1--3 to $\times$0.8--1.5 recovers most of the metal-off
gain (teeth +0.024) with a smaller numbering cost. Widening the image-level ranges (photon count, read noise, detector
blur, background subtraction, unsharp masking, gamma) improves teeth Dice by 0.031 and FDI Dice by 0.033 in every seed,
with numbering unchanged. It is the only change that improves outline and per-tooth Dice at no cost to enumeration.

\begin{table}[t]
\centering\small
\caption{Follow-up experiments on the randomization ranges (DENTEX test, three seeds, ablation schedule).}
\label{tab:follow}
\begin{adjustbox}{max width=\linewidth}
\begin{tabular}{@{}lccccc@{}}
\toprule
Variant & Teeth Dice & FDI Dice & Numbering & $\Delta$ teeth [min, max] & $\Delta$ numb. \\
\midrule
Reference & \pmsd{0.734}{0.004} & \pmsd{0.569}{0.015} & \pmsd{0.911}{0.004} & -- & -- \\
Metal off & \pmsd{0.767}{0.009} & \pmsd{0.592}{0.007} & \pmsd{0.895}{0.015} & +0.033 [+0.022, +0.047] & $-$0.015 \\
Metal $\times$0.8--1.5 & \pmsd{0.758}{0.005} & \pmsd{0.584}{0.002} & \pmsd{0.901}{0.006} & +0.024 [+0.014, +0.031] & $-$0.009 \\
Wider image ranges & \pmsd{0.766}{0.011} & \pmsd{0.601}{0.019} & \pmsd{0.914}{0.009} & +0.031 [+0.020, +0.042] & +0.003 \\
Both & \pmsd{0.779}{0.009} & \pmsd{0.599}{0.007} & \pmsd{0.894}{0.012} & +0.044 [+0.036, +0.058] & $-$0.017 \\
\bottomrule
\end{tabular}
\end{adjustbox}
\end{table}

We then retrained the full model with the narrower metal range and the wider image ranges, three seeds each, and
compared it with three seeds of the original recipe (Table~\ref{tab:retune}). On real panoramics the retuned model is
better and more stable: teeth Dice \pmsd{0.773}{0.008} vs.\ \pmsd{0.755}{0.020}, and detection 0.914 vs.\ 0.871, with
unchanged numbering. But it is worse on the structures that lack real labels. On synthetic test volumes the IAC falls from
0.746 / 0.755 to 0.721 / 0.721 and implants from 0.687 to 0.573. The canal loss carries over to real images: a crop model
trained with the retuned recipe reaches a CONTACT canal Dice of \pmsd{0.397}{0.084}, against \pmsd{0.479}{0.039} for the
original, and it is lower in every seed pair (Table~\ref{tab:crop}). Stronger image-level randomization thus trades the
thin, low-contrast canal against the teeth. We keep the original recipe as our main model. The ablation's conclusions are
drawn on teeth and should not be generalised to every structure.

\begin{table}[t]
\centering\small
\caption{Full-schedule zero-shot models with the original and the retuned randomization ranges, mean $\pm$ sd over three
seeds each (DENTEX test and synthetic test volumes).}
\label{tab:retune}
\begin{adjustbox}{max width=\linewidth}
\begin{tabular}{@{}lcccc@{}}
\toprule
& \multicolumn{2}{c}{Original ranges} & \multicolumn{2}{c}{Retuned ranges} \\
\midrule
DENTEX teeth Dice & \multicolumn{2}{c}{\pmsd{0.755}{0.020}} & \multicolumn{2}{c}{\pmsd{0.773}{0.008}} \\
DENTEX FDI Dice & \multicolumn{2}{c}{\pmsd{0.582}{0.019}} & \multicolumn{2}{c}{\pmsd{0.598}{0.013}} \\
DENTEX detection & \multicolumn{2}{c}{\pmsd{0.871}{0.038}} & \multicolumn{2}{c}{\pmsd{0.914}{0.014}} \\
DENTEX numbering & \multicolumn{2}{c}{\pmsd{0.900}{0.001}} & \multicolumn{2}{c}{\pmsd{0.895}{0.016}} \\
Synthetic IAC Dice, left / right & \multicolumn{2}{c}{0.746 / 0.755} & \multicolumn{2}{c}{0.721 / 0.721} \\
Synthetic pulp Dice & \multicolumn{2}{c}{0.669} & \multicolumn{2}{c}{0.655} \\
Synthetic implant Dice & \multicolumn{2}{c}{0.687} & \multicolumn{2}{c}{0.573} \\
\bottomrule
\end{tabular}
\end{adjustbox}
\end{table}

\subsection{Randomization versus no randomization and image translation}
Without randomization, simulation does not transfer at all (Table~\ref{tab:sources}). Trained on nominal renders only,
the model reaches a teeth Dice of \pmsd{0.277}{0.070} on DENTEX, detects 8\% of the teeth and numbers 19\% of those
correctly, even though it fits the synthetic data well. This mirrors the ablation, where removing image-level noise alone
cost 0.44 Dice.

The CycleGAN-translated renders look convincingly real (Fig.~\ref{fig:cyclegan}), yet segmentation trained on them also
transfers far worse than physics-based randomization: teeth Dice 0.562 vs.\ 0.755, and numbering accuracy 0.104 vs.\
0.900. Adding our image-level randomization before translation does not help (0.509 / 0.216). The samples show why: the
unpaired translation changes the anatomy. Bright restorations disappear, and tooth outlines and gaps shift, so the
projected labels no longer describe the image. Per-tooth identity suffers most, because it depends on exactly the details
the translation alters. A simulator whose labels stay correct under every randomization avoids this failure. One caveat:
this is a single run of the standard architecture, without structure-preserving losses.

\begin{table}[t]
\centering\small
\caption{Zero-shot results by source of training images (DENTEX test).}
\label{tab:sources}
\begin{adjustbox}{max width=\linewidth}
\begin{tabular}{@{}lccccc@{}}
\toprule
Training images & Runs & Teeth Dice & FDI Dice & Detection & Numbering \\
\midrule
Nominal renders, no randomization & 3 & \pmsd{0.277}{0.070} & \pmsd{0.013}{0.003} & \pmsd{0.079}{0.019} & \pmsd{0.190}{0.138} \\
Physics-based randomization (ours) & 3 & \textbf{\pmsd{0.755}{0.020}} & \textbf{\pmsd{0.582}{0.019}} & \textbf{\pmsd{0.871}{0.038}} & \textbf{\pmsd{0.900}{0.001}} \\
CycleGAN-translated renders & 1 & 0.562 & 0.125 & 0.463 & 0.104 \\
Randomization, then CycleGAN & 1 & 0.509 & 0.145 & 0.344 & 0.216 \\
\bottomrule
\end{tabular}
\end{adjustbox}
\end{table}

\begin{figure}[t]
\centering
\includegraphics[width=\textwidth]{fig_cyclegan_samples_DATA_paper.png}
\caption{CycleGAN translation between nominal synthetic renders and real DENTEX panoramics~\cite{hamamci2023}. The
translated images look realistic, but the translation alters anatomy: restorations vanish or appear and tooth outlines
shift, so the projected labels no longer match the translated image.}
\label{fig:cyclegan}
\end{figure}

\subsection{Third molar, canal and contact on real crops}
\label{sec:contact}

\textbf{Segmentation.} The zero-shot crop model segments the third molar on the 444 CONTACT test crops with a Dice of
0.611 and the canal with 0.436 (seed 0, which initialises the contact experiments). Crop models vary considerably with
the training seed. Over three seeds, molar Dice is \pmsd{0.672}{0.058} and canal Dice \pmsd{0.479}{0.039}, and seed 0
is the weakest (Table~\ref{tab:crop}). Molar outlines are often tight, with per-image Dice of 0.85--0.95
(Fig.~\ref{fig:contact}). Failures mostly come from segmenting the neighbouring second molar. CONTACT's canal masks are
drawn as wide bands between the cortical lines, while our projected label is the canal lumen, which caps canal Dice.

\begin{table}[t]
\centering\small
\caption{Zero-shot crop models on the CONTACT test set (matched resolution), three seeds each.}
\label{tab:crop}
\begin{adjustbox}{max width=\linewidth}
\begin{tabular}{@{}lcccc@{}}
\toprule
Crop model & Molar Dice & Canal Dice & Canal Dice, group 3 & Canal Dice per seed \\
\midrule
Original ranges & \pmsd{0.672}{0.058} & \pmsd{0.479}{0.039} & \pmsd{0.490}{0.036} & 0.436, 0.510, 0.492 \\
Retuned ranges & \pmsd{0.662}{0.085} & \pmsd{0.397}{0.084} & \pmsd{0.399}{0.069} & 0.424, 0.465, 0.303 \\
\bottomrule
\end{tabular}
\end{adjustbox}
\end{table}

\begin{figure}[t]
\centering
\includegraphics[width=\textwidth]{fig_contact_overlays_DATA_paper.png}
\caption{Zero-shot crop-model predictions on CONTACT test crops~\cite{ulusoy2025} (group 1, matched resolution).
CONTACT's canal masks are wide bands between the cortical lines; our projected canal label is the narrower lumen.}
\label{fig:contact}
\end{figure}

\textbf{2D overlap cannot detect contact.} The predicted molar--canal overlap fraction gives a contact AUC of 0.510, and
0.468 for contact vs.\ overlap without contact (group 1 vs.\ 3). Both are at chance, which is the radiological problem
CONTACT was built to study.

\textbf{The contact head learns a weak but real signal from CBCT alone.} Trained only on synthetic crops with labels from
3D geometry, the head reaches an AUC of 0.596 [0.538, 0.656] on the real test set, and 0.575 for group 1 vs.\ 3
(Table~\ref{tab:contact}). For reference, the six expert-scored radiographic signs supplied with CONTACT give an AUC of
0.54 when summed, and 0.72 for the best single sign (interruption of the canal's white line). Looser contact definitions
transfer worse ($\leq$\,1\,mm: 0.564; $\leq$\,2\,mm: 0.515), and the validation set agrees, so we keep the a-priori
definition. Validation AUC on CONTACT (0.77--0.78) is well above test AUC. The official test set appears intrinsically
harder: even the ground-truth mask overlap separates contact with an AUC of 0.717 on test against 0.763 on validation, and
the test crops are blurrier.

\begin{table}[t]
\centering\small
\caption{Zero-shot contact prediction on the CONTACT test set (444 crops, matched resolution).}
\label{tab:contact}
\begin{adjustbox}{max width=\linewidth}
\begin{tabular}{@{}lcccc@{}}
\toprule
Model (zero-shot) & Molar Dice & Canal Dice & Contact AUC [95\% CI] & Group 1 vs.\ 3 AUC \\
\midrule
Crop model, 2D overlap readout & 0.611 & 0.436 & 0.510 & 0.468 \\
Contact head, contact $\leq$\,0.52\,mm & 0.604 & 0.478 & \textbf{0.596 [0.538, 0.656]} & \textbf{0.575} \\
Contact head, contact $\leq$\,1.0\,mm & 0.624 & 0.412 & 0.564 [0.508, 0.622] & 0.554 \\
Contact head, contact $\leq$\,2.0\,mm & 0.536 & 0.453 & 0.515 [0.457, 0.576] & 0.521 \\
\bottomrule
\end{tabular}
\end{adjustbox}
\end{table}

\textbf{Few-shot contact.} CBCT pretraining helps when real labels are scarce, and real-only training is better once
they are plentiful (Table~\ref{tab:fewcontact}, Fig.~\ref{fig:fewcontact}). With nine labelled crops, \sr\ reaches a
contact AUC of \pmsd{0.579}{0.017} against \pmsd{0.445}{0.043}, and a canal Dice of 0.674 vs.\ 0.495. The real-only
model at this budget is below chance, so it has learned a spurious cue. With 44 crops \sr\ still has the higher contact
AUC (\pmsd{0.616}{0.009} vs.\ \pmsd{0.589}{0.012}), with equal canal Dice. At 88 and 220 crops the contact AUC of the
two arms is level, but real-only training segments the canal better. From 439 crops real-only training is ahead on every
measure, reaching a contact AUC of \pmsd{0.745}{0.030} with all 878 crops against \pmsd{0.698}{0.012}. The synthetic
prior does not carry over to CONTACT's harder test split once enough real labels are available. A plausible cause for
the canal gap is the label definition: our labels mark the canal lumen, CONTACT's masks a wider band. A fine-tuned model
has to unlearn the narrower definition, and the higher learning rate, which lets it move further from its
initialisation, narrows the gap (with all crops, canal Dice 0.846 at learning rate $10^{-3}$ vs.\ 0.829 at
$3\cdot10^{-4}$).

\begin{figure}[t]
\centering
\includegraphics[width=\textwidth]{fig_fewshot_contact.pdf}
\caption{Few-shot contact on the CONTACT test set as a function of the number of labelled real crops (mean over three
seeds; bands show $\pm$ one sd). Both arms use the same schedule. The dashed line is the zero-shot contact model; the
dotted line is chance.}
\label{fig:fewcontact}
\end{figure}

\begin{table}[t]
\centering\small
\caption{Few-shot contact on the CONTACT test set, mean $\pm$ sd over three seeds. Both arms are trained with learning
rate $10^{-3}$ for 90 epochs. RO = real-only (ImageNet initialisation); S2R = fine-tuned contact model. Bold: the better
arm where the means differ by more than one sd of either arm.}
\label{tab:fewcontact}
\begin{adjustbox}{max width=\linewidth}
\begin{tabular}{@{}rlccc@{}}
\toprule
Labelled crops & Arm & Contact AUC & Group 1 vs.\ 3 AUC & Canal Dice \\
\midrule
\multirow{2}{*}{9 (1\%)} & RO & \pmsd{0.445}{0.043} & \pmsd{0.446}{0.033} & \pmsd{0.495}{0.064} \\
 & S2R & \textbf{\pmsd{0.579}{0.017}} & \textbf{\pmsd{0.529}{0.016}} & \textbf{\pmsd{0.674}{0.024}} \\
\midrule
\multirow{2}{*}{44 (5\%)} & RO & \pmsd{0.589}{0.012} & \pmsd{0.579}{0.027} & \pmsd{0.760}{0.029} \\
 & S2R & \textbf{\pmsd{0.616}{0.009}} & \pmsd{0.570}{0.010} & \pmsd{0.776}{0.016} \\
\midrule
\multirow{2}{*}{88 (10\%)} & RO & \pmsd{0.594}{0.089} & \pmsd{0.559}{0.103} & \textbf{\pmsd{0.811}{0.008}} \\
 & S2R & \pmsd{0.633}{0.020} & \pmsd{0.571}{0.022} & \pmsd{0.794}{0.005} \\
\midrule
\multirow{2}{*}{220 (25\%)} & RO & \pmsd{0.658}{0.028} & \pmsd{0.597}{0.034} & \textbf{\pmsd{0.839}{0.007}} \\
 & S2R & \pmsd{0.654}{0.004} & \pmsd{0.587}{0.003} & \pmsd{0.817}{0.006} \\
\midrule
\multirow{2}{*}{439 (50\%)} & RO & \textbf{\pmsd{0.722}{0.027}} & \textbf{\pmsd{0.683}{0.034}} & \textbf{\pmsd{0.858}{0.004}} \\
 & S2R & \pmsd{0.690}{0.002} & \pmsd{0.632}{0.005} & \pmsd{0.828}{0.002} \\
\midrule
\multirow{2}{*}{878 (100\%)} & RO & \textbf{\pmsd{0.745}{0.030}} & \textbf{\pmsd{0.697}{0.038}} & \textbf{\pmsd{0.866}{0.006}} \\
 & S2R & \pmsd{0.698}{0.012} & \pmsd{0.639}{0.012} & \pmsd{0.846}{0.009} \\
\bottomrule
\end{tabular}
\end{adjustbox}
\end{table}

\begin{table}[t]
\centering\small
\caption{Sensitivity of the \sr\ arm to its fine-tuning schedule (test set, mean over three seeds; columns give the
fraction of the real training pool). The reported schedules were chosen without the test set
(Section~\ref{sec:setup}).}
\label{tab:sched}
\setlength{\tabcolsep}{4pt}
\begin{adjustbox}{max width=\linewidth}
\begin{tabular}{@{}llcccccc@{}}
\toprule
Data, metric & \Sr\ schedule & 1\% & 5\% & 10\% & 25\% & 50\% & 100\% \\
\midrule
DENTEX, FDI Dice & lr $3\cdot10^{-4}$, 40 epochs (reported) & 0.702 & 0.720 & 0.736 & 0.749 & 0.762 & 0.774 \\
 & lr $3\cdot10^{-4}$, 120 epochs & 0.699 & 0.715 & 0.730 & 0.746 & 0.760 & 0.775 \\
\midrule
CONTACT, contact AUC & lr $3\cdot10^{-4}$, 30 epochs & 0.582 & 0.601 & 0.621 & 0.627 & 0.649 & 0.672 \\
 & lr $3\cdot10^{-4}$, 90 epochs & 0.582 & 0.603 & 0.622 & 0.633 & 0.655 & 0.668 \\
 & lr $10^{-3}$, 90 epochs (reported) & 0.579 & 0.616 & 0.633 & 0.654 & 0.690 & 0.698 \\
\bottomrule
\end{tabular}
\end{adjustbox}
\end{table}

% =====================================================================================================================
\section{Discussion}

\textbf{What simulation buys.} The clearest result is tooth enumeration. A model that has only seen simulated panoramics
numbers 90\% of detected real teeth correctly. Four real images raise this to 93\%, a level that real-only training
reaches only with about 40 or more labelled images. Two features of the simulated data probably explain this. First,
every render has a complete, exact FDI label for every tooth in the volume, including teeth that are hard to see because
they overlap or are blurred. Second, the projection preserves the 3D order of the arch. Real-only training with few
labels learns the outline of the dentition but not which tooth is which, and that is the information simulation
supplies. The benefit is limited to the low-data regime: once about 100 labelled images are available, a well-tuned
real-only model does as well. For teeth, which already have public labels, simulation is therefore mainly a way to need
far fewer of them. Its larger promise is for structures without any public labels, which our evaluation can only partly
test.

\textbf{What makes a panoramic simulator work.} The ablation suggests a division of labour. The image-level factors
(noise, detector blur and vendor processing) decide whether anything transfers at all. Without noise, real teeth Dice
collapses to about 0.30 even though synthetic accuracy is unchanged. The geometry factors, which are what make a
panoramic distinct, matter less for the outline but consistently help tooth identity. Our reading is that blur, ghosts
and trough variation create the confusions between neighbouring teeth that occur on real images, and the model learns to
resolve them. Metal randomization at its original strength made restorations unrealistically dominant and hurt the
outline. These per-factor effects are small (a few Dice points), and we needed three seeds and paired renders to separate
them from noise. Unpaired ablations with a single seed would have been misleading here.

\textbf{Teeth are not the whole story.} Retuning the ranges for teeth hurt the canal, on synthetic and real images alike.
Different structures want different randomization: strong noise and processing wash out a thin, low-contrast canal that
a tooth survives easily. Simulation studies that tune or ablate only on the structure with real labels may be optimising
the wrong objective for the rest. Validating each structure, or randomizing per structure, is a natural next step.

\textbf{Contact from geometry.} A panoramic cannot show 3D contact directly, and the predicted 2D overlap indeed carries
no information. A head trained on contact labels computed from CBCT recovers a weak but significant signal (AUC 0.60,
confidence interval above 0.5) from the image alone, and it does so without a single real contact label. It is not
clinically useful at this level. Its value is that it shows the 3D relation leaves a learnable trace in the projection,
and that this trace can be learned from simulation. As a pretraining prior it pays off up to about 40 real crops; at
around 100--200 crops the arms are level, and from about 400 a real-only model is clearly better. Our contact definition
and CONTACT's may not coincide: we use labels touching on the canal lumen, CONTACT uses a CBCT-confirmed relation.
CONTACT also counts far more contact (74\%) than our synthetic population (33\%). The threshold analysis rules out the
simplest explanation, a looser threshold, but other differences in definition may remain.

\textbf{Translation versus simulation.} CycleGAN produced the most realistic-looking images and the worst segmentation.
For dense-label tasks this argues against bridging the sim-to-real gap by changing the image after the fact. Changing it
inside the simulator keeps the labels correct by construction.

\textbf{Limitations.} Real-image evaluation covers teeth (DENTEX) and the third molar, canal and contact (CONTACT). The
sinus, pulp and incisive canal are evaluated on synthetic data only; we found no public expert panoramic labels for them,
and we did not produce our own. DENTEX tooth polygons are coarse, and Dice rewards learning their style, which favours
models trained on DENTEX labels. CONTACT is distributed as crops without patient identifiers, so our validation split may
share patients with training, and its official test set appears harder than its training pool. Most results use three
seeds, and the seed-to-seed spread is not negligible. The panoramic model we report and use for initialisation is the
best of three seeds on DENTEX (0.774 against a mean of 0.755), and the crop model we use is the weakest of three on
CONTACT. The few-shot and contact experiments vary the fine-tuning seed but start from this one model. Training schedules
were tuned per arm, not per label budget; a per-budget search might shift the crossover points somewhat, and on CONTACT
the \sr\ result depended on the learning rate. The CycleGAN baseline is a single run of the standard architecture.
Finally, ToothFairy3's field of view limits the simulation: most volumes do not include the upper jaw, orbits or
condyles, the patients bite in an open position rather than on a bite block, and the simulated resolution (0.3\,mm) is
coarser than that of a real panoramic.

% =====================================================================================================================
\section{Conclusion}

A physics-based panoramic simulator over labelled CBCT, with randomization of the panoramic geometry and image formation,
is enough to train a model that segments and numbers teeth on real panoramics without any real labels. With a handful of
real images, it matches real-only training on about ten times as many labels, although the advantage disappears once
about 100 labelled images are available. Without randomization, simulation does not transfer. The paired ablation shows
which parts of the simulation matter, and that the answer differs between structures. Simulated labels also carry
information that no 2D annotation can: a contact head trained on 3D geometry alone detects true third-molar--canal
contact above chance, where the visible overlap cannot, although a few hundred real contact labels overtake it.

% =====================================================================================================================
\section*{CRediT authorship contribution statement}
\textbf{Dipesh Tripathi:} Conceptualization, Methodology, Software, Validation, Formal analysis, Investigation, Data
curation, Writing -- original draft, Writing -- review \& editing, Visualization.

\section*{Declaration of competing interest}
The author declares no known competing financial interests or personal relationships that could have
appeared to influence the work reported in this paper.

\section*{Funding}
This research did not receive any specific grant from funding agencies in the public, commercial, or not-for-profit
sectors.

\section*{Data availability}
ToothFairy3, DENTEX and CONTACT are publicly available from their providers under their respective licences. The
simulator source code, configuration files, data splits and trained models are available from the author on reasonable
request.

\section*{Acknowledgements}
Computations were performed on the Lawrence high-performance computing cluster at the University of South Dakota.

\bibliographystyle{elsarticle-num}
\bibliography{refs}

\end{document}
